%============================================================
% MTH 112Z Project - Supplement Angles
% Updated 202302
%============================================================

\section{Angles} \label{supplement-angles}
In this Section, we will learn some things.

\subsection*{Coterminal Angles} \label{supplement-coterminal-angles}
  \begin{myDefinition}
    Two angles are \textbf{coterminal} if they have the same terminal side when in standard position.
  \end{myDefinition}

      Since \ang{360} represents a complete revolution, if we add integer-multiples of \ang{360} to an angle measured in degrees, we'll obtain a coterminal angle. Similarly, since $2\pi$ represents a complete revolution in radians, if we add integer-multiples of $2\pi$ to an angle measured in radians, we'll obtain a coterminal angle. We can summarize this information as follows

      If $\theta$ is measured in degrees, $\theta$ and $\theta + \ang{360} \cdot k$, where $k\in\mathbb{Z}$, are coterminal.
    
      If $\theta$ is measured in radians, $\theta$ and $\theta + 2\pi \cdot k$, where $k\in\mathbb{Z}$, are coterminal.
    
      \begin{myExample}
        The angles $\ang{45}$, $\ang{405}$, and $\ang{-315}$ are coterminal as illustrated inFigure \ref{fig:coterminal}.
      \end{myExample}
      
      \begin{center}
        \captionsetup{type=figure} \captionof{figure}{Coterminal angles}
          %<description>This is a graph showing coterminal angles of 45 degrees, 405 degrees, and -315 degrees which all have the same terminal side.</description>
          
            \begin{tikzpicture}
              \begin{axis}[
                  width=8cm,
                  grid=none,
                  ticks=none,
                  trig format plots=rad,
                  axis equal,
                  xmin=,
                  xmax=,
                  ymin=,
                  ymax=,
                ]
                \addplot[ray] coordinates {(0,0) (11,11)};
                \addplot[angle, domain=0:pi/4] ({7*cos(x)}, {7*sin(x)}) node[pos=0.5, above right, inner sep=1pt] {$45^{\circ}$};
                \addplot[angle, domain=0:-7*pi/4, samples=50] ({10*cos(x)}, {10*sin(x)}) node[pos=0.75, above left, inner sep=0pt] {$-315^{\circ}$};
                \addplot[angle, domain=0:9*pi/4, samples=50] ({4*log10(x+10)*cos(x)}, {4*log10(x+10)*sin(x)}) node[pos=0.5, below left, inner sep=0pt] {$405^{\circ}$};
              \end{axis}
            \end{tikzpicture}\label{fig:coterminal}
      \end{center}
    </example>

    \subsection*{Reference Angles}
    \begin{myDefinition}
        
          The \textbf{reference angle} for an angle in standard position is the positive acute angle formed by the $x$-axis and the terminal side of the angle.
        
    \end{myDefinition}
    
      Depending on the location of the angle's terminal side, we'll have to use a different calculation to determine the angle's reference angle.
    
    <example>
      
        The angles $\frac{\pi}{3}$ and $30^{\circ}$ are their own reference angles since they are acute angles; seen inFigure \ref{fig:example-ref-pi3} andFigure \ref{fig:example-ref-30}.
      

      <sidebyside widths="37.5% 37.5%">
        \begin{center}
          \captionsetup{type=figure} \captionof{figure}{}
            %<description>Graph showing the angle pi/3 radians.</description>
              \begin{tikzpicture}
                \begin{axis}[
                    width=6cm,
                    grid=none,
                    ticks=none,
                    trig format plots=rad,
                    axis equal,
                  ]
                  \addplot[ray] coordinates {(0,0) (3,5.196)};
                  \addplot[angle, domain=0:pi/3] ({3.5*cos(x)}, {3.5*sin(x)}) node[pos=0.5, above right, inner sep=0pt] {$\frac{\pi}{3}$};
                \end{axis}
              \end{tikzpicture}\label{fig:example-ref-pi3}
        \end{center}

        \begin{center}
          \captionsetup{type=figure} \captionof{figure}{}
            %<description>Graph showing the angle 30 degrees.</description>
              \begin{tikzpicture}
                \begin{axis}[
                    width=6cm,
                    grid=none,
                    ticks=none,
                    trig format plots=rad,
                    axis equal,
                  ]
                  \addplot[ray] coordinates {(0,0) (5.196,3)};
                  \addplot[angle, domain=0:pi/6] ({3.5*cos(x)}, {3.5*sin(x)}) node[pos=0.6, right] {$30^{\circ}$};
                \end{axis}
              \end{tikzpicture}\label{fig:example-ref-30}
        \end{center}
      </sidebyside>
    </example>
    <example>
      
        The reference angle for $\frac{2\pi}{3}$ is $\pi-\frac{2\pi}{3}=\frac{\pi}{3}$ (seeFigure \ref{fig:example-ref-2pi3}), while the reference angle for $150^{\circ}$ is $180^{\circ}-150^{\circ}=30^{\circ}$ (seeFigure \ref{fig:example-ref-150}).
      
      <sidebyside widths="37.5% 37.5%">
        \begin{center}
          \captionsetup{type=figure} \captionof{figure}{}
            %<description>Graph showing the angle (2pi)/3 radians which has a reference angle of pi/3 radians.</description>
              \begin{tikzpicture}
                \begin{axis}[
                    width=6cm,
                    grid=none,
                    ticks=none,
                    trig format plots=rad,
                    axis equal,
                  ]
                  \addplot[ray] coordinates {(0,0) (-3,5.196)};
                  \addplot[angle, domain=0:2*pi/3] ({3.5*cos(x)}, {3.5*sin(x)}) node[pos=0.5, above right, inner sep=0pt] {$\frac{2\pi}{3}$};
                  \addplot[angle, domain=pi:2*pi/3]({2.5*cos(x)}, {2.5*sin(x)}) node[pos=0.6, left] {$\frac{\pi}{3}$};
                \end{axis}
              \end{tikzpicture}\label{fig:example-ref-2pi3}
        \end{center}

        \begin{center}
          \captionsetup{type=figure} \captionof{figure}{}
            %<description>Graph showing the angle 150 degrees which has a reference angle of 30 degrees.</description>
              \begin{tikzpicture}
                \begin{axis}[
                    width=6cm,
                    grid=none,
                    ticks=none,
                    trig format plots=rad,
                    axis equal,
                  ]
                  \addplot[ray] coordinates {(0,0) (-5.196,3)};
                  \addplot[angle, domain=0:5*pi/6] ({3.5*cos(x)}, {3.5*sin(x)}) node[pos=0.5, above right, inner sep=1pt] {$150^{\circ}$};
                  \addplot[angle, domain=pi:5*pi/6]({2.5*cos(x)}, {2.5*sin(x)}) node[pos=0.6, left] {$30^{\circ}$};
                \end{axis}
              \end{tikzpicture}\label{fig:example-ref-150}
        \end{center}
      </sidebyside>
    </example>

    <example>
      
        The reference angle for $\frac{4\pi}{3}$ is $\frac{4\pi}{3}-\pi=\frac{\pi}{3}$ (seeFigure \ref{fig:example-ref-4pi3}), while the reference angle for $210^{\circ}$ is $210^{\circ}-180^{\circ}=30^{\circ}$ (seeFigure \ref{fig:example-ref-210}).
      
      <sidebyside widths="37.5% 37.5%">
        \begin{center}
          \captionsetup{type=figure} \captionof{figure}{}
            %<description>Graph showing the angle (4pi)/3 radians which has a reference angle of pi/3 radians.</description>
              \begin{tikzpicture}
                \begin{axis}[
                    width=6cm,
                    grid=none,
                    ticks=none,
                    trig format plots=rad,
                    axis equal,
                  ]
                  \addplot[ray] coordinates {(0,0) (-3,-5.196)};
                  \addplot[angle, domain=0:4*pi/3] ({2.5*cos(x)}, {2.5*sin(x)}) node[pos=0.5, above left, inner sep=0pt] {$\frac{4\pi}{3}$};
                  \addplot[angle, domain=pi:4*pi/3]({3.5*cos(x)}, {3.5*sin(x)}) node[pos=0.5, below left, inner sep=0pt] {$\frac{\pi}{3}$};
                \end{axis}
              \end{tikzpicture}\label{fig:example-ref-4pi3}
        \end{center}

        \begin{center}
          \captionsetup{type=figure} \captionof{figure}{}
            %<description>Graph showing the angle 210 degrees which has a reference angle of 30 degrees.</description>
              \begin{tikzpicture}
                \begin{axis}[
                    width=6cm,
                    grid=none,
                    ticks=none,
                    trig format plots=rad,
                    axis equal,
                  ]
                  \addplot[ray] coordinates {(0,0) (-5.196,-3)};
                  \addplot[angle, domain=0:7*pi/6] ({2.5*cos(x)}, {2.5*sin(x)}) node[pos=0.5, above left, inner sep=0pt] {$210^{\circ}$};
                  \addplot[angle, domain=pi:7*pi/6]({3.5*cos(x)}, {3.5*sin(x)}) node[pos=0.6, left] {$30^{\circ}$};
                \end{axis}
              \end{tikzpicture}\label{fig:example-ref-210}
        \end{center}
      </sidebyside>
    </example>

    <example>
      
        The reference angle for $\frac{5\pi}{3}$ is $2\pi-\frac{5\pi}{3}=\frac{\pi}{3}$ (seeFigure \ref{fig:example-ref-5pi3}), while the reference angle for $330^{\circ}$ is $\ang{360}-330^{\circ}=30^{\circ}$ (seeFigure \ref{fig:example-ref-330}).
      
      <sidebyside widths="37.5% 37.5%">
        \begin{center}
          \captionsetup{type=figure} \captionof{figure}{}
            %<description>Graph showing the angle (5pi)/3 radians which has a reference angle of pi/3 radians.</description>
              \begin{tikzpicture}
                \begin{axis}[
                    width=6cm,
                    grid=none,
                    ticks=none,
                    trig format plots=rad,
                    axis equal,
                  ]
                  \addplot[ray] coordinates {(0,0) (3,-5.196)};
                  \addplot[angle, domain=0:5*pi/3, samples=50] ({2.5*cos(x)}, {2.5*sin(x)}) node[pos=0.5, above left, inner sep=0pt] {$\frac{5\pi}{3}$};
                  \addplot[angle, domain=0:-pi/3]({3.5*cos(x)}, {3.5*sin(x)}) node[pos=0.6, right] {$\frac{\pi}{3}$};
                \end{axis}
              \end{tikzpicture}\label{fig:example-ref-5pi3}
        \end{center}

        \begin{center}
          \captionsetup{type=figure} \captionof{figure}{}
            %<description>Graph showing the angle 330 degrees which has a reference angle of 30 degrees.</description>
              \begin{tikzpicture}
                \begin{axis}[
                    width=6cm,
                    grid=none,
                    ticks=none,
                    trig format plots=rad,
                    axis equal,
                  ]
                  \addplot[ray] coordinates {(0,0) (5.196,-3)};
                  \addplot[angle, domain=0:11*pi/6, samples=50] ({2.5*cos(x)}, {2.5*sin(x)}) node[pos=0.5, above left, inner sep=0pt] {$330^{\circ}$};
                  \addplot[angle, domain=0:-pi/6]({3.5*cos(x)}, {3.5*sin(x)}) node[pos=0.6, right] {$30^{\circ}$};
                \end{axis}
              \end{tikzpicture}\label{fig:example-ref-330}
        \end{center}
      </sidebyside>
    </example>

    <example>
      
        The reference angle for $7.5$ radians is $7.5-2\pi\approx 1.2$ radians (seeFigure \ref{fig:example-ref-7point5}), and the reference angle for $-137^{\circ}$ is $180^{\circ}+( \, -137^{\circ} ) \,=43^{\circ}$ (seeFigure \ref{fig:example-ref-negative137}).
      
      <sidebyside widths="37.5% 37.5%">
        \begin{center}
          \captionsetup{type=figure} \captionof{figure}{}
            %<description>Graph showing the angle 7.5 radians which has a reference angle of 1.2 radians.</description>
              \begin{tikzpicture}
                \begin{axis}[
                    width=6cm,
                    grid=none,
                    ticks=none,
                    trig format plots=rad,
                    axis equal,
                  ]
                  \addplot[ray] coordinates {(0,0) (2.216,6)};
                  \addplot[angle, domain=0:7.5, samples=50] ({0.7*ln(x+10)*cos(x)}, {0.7*ln(x+10)*sin(x)}) node[pos=0.5, below left, inner sep=1pt] {$7.5$};
                  \addplot[angle, domain=0:7.5-2*pi]({3*cos(x)}, {3*sin(x)}) node[pos=0.6, above right, inner sep=1pt] {$7.5-2\pi$};
                \end{axis}
              \end{tikzpicture}\label{fig:example-ref-7point5}
        \end{center}

        \begin{center}
          \captionsetup{type=figure} \captionof{figure}{}
            %<description>Graph showing the angle -137 degrees which has a reference angle of 43 degrees.</description>
              \begin{tikzpicture}
                \begin{axis}[
                    width=6cm,
                    grid=none,
                    ticks=none,
                    axis equal,
                  ]
                  \addplot[ray] coordinates {(0,0) (-4.388,-4.092)};
                  \addplot[angle, domain=0:-137] ({3.5*cos(x)}, {3.5*sin(x)}) node[pos=0.5, below right, inner sep=0pt] {$-137^{\circ}$};
                  \addplot[angle, domain=180:223]({2.5*cos(x)}, {2.5*sin(x)}) node[pos=0.6, left] {$43^{\circ}$};
                \end{axis}
              \end{tikzpicture}\label{fig:example-ref-negative137}
        \end{center}
      </sidebyside>
    </example>

  </subsection>

  
    \subsection*{Degrees, Minutes, and Seconds}
    
      When measuring angles in degrees, fractions of a degree can be represented in minutes and seconds.
    
    \begin{myDefinition}
        
          One \textbf{minute} is $\frac{1}{60}$ of a degree, so $60$ minutes is equal to one degree (written $
          60'=1^{\circ}$).
        
        
          One \textbf{second} is $\frac{1}{60}$ of a minute, i.e., $\frac{1}{3600}$ of a degree, so $60$ seconds is equal to one minute (written $60''=1'$) and $3600$ seconds is equal to one degree (written $3600''=1^{\circ}$).
        
    \end{myDefinition}
    
      When we write an angle's measure in the form $D^{\circ}M'S''$ (where $D$, $M$, and $S$ are real numbers), then the angle's measure is $D$ degrees plus $M$ minutes plus $S$ seconds.
    
    <example>
        
          Convert $34^{\circ}15'27''$ into decimal form.
        
      <solution>
          \begin{align*}
            \ang{34;15;27} &= \ang{34}+15'\left( \frac{1^{\circ}}{60'} \right)+27'' \left(\frac{1^{\circ}}{3600''}\right) \\
             &= 34^{\circ}+0.25^{\circ}+0.0075^{\circ} \\
             &= 34.2575^{\circ} \\
          \end{align*}
        
      </solution>
    </example>
    <example>
        
          Convert $61.72^{\circ}$ into $D^{\circ}M'S''$ form.
        
      <solution>
        
          \begin{align*}
             61.72^{\circ} &= 61^{\circ}+0.72^{\circ} \\
             &= 61^{\circ}+0.72^{\circ} \cdot \left( \frac{60'}{1^{\circ}} \right) \\
             &= 61^{\circ}+43.2' \\
             &= 61^{\circ}43'+0.2' \\
             &= 61^{\circ}43'+0.2' \cdot \left( \frac{60''}{1'} \right) \\
             &= 61^{\circ}43'12''\\
          \end{align*}
        
      </solution>
    </example>
  </subsection>

    <exercises xml:id="exercises-angles">

    <exercisegroup cols="3" xml:id="exercisegroup-coterminal">
      \subsection*{Coterminal Angles}
        
            In <xref ref="exercisegroup-coterminal" text="local">Exercises</xref>, find both a positive and negative angle that is coterminal angle with the following angles.
        
        <exercise>
                $63^{\circ}$
            <answer>
               $423^{\circ}$ and $-297^{\circ}$ are coterminal with $63^{\circ}$.
            </answer>
        </exercise>
        <exercise>
                $\frac{\pi}{9}$
                <answer>
                    $\frac{19\pi}{9}$ and $-\frac{17\pi}{9}$ are coterminal with $\frac{\pi}{9}$.
                </answer>
        </exercise>
        <exercise>
                $\frac{13\pi}{8}$
                <answer>
                    $\frac{29\pi}{8}$ and $-\frac{3\pi}{8}$ are coterminal with $\frac{13\pi}{8}$.
                </answer>
        </exercise>
    </exercisegroup>
    <exercisegroup cols="3" xml:id="exercisegroup-reference">
      \subsection*{Reference Angles}
        
            In <xref ref="exercisegroup-reference" text="local">Exercises</xref>, find the reference angle for the following angles.
        
        <exercise>
                $120^{\circ}$
            <answer>
               $60^{\circ}$
            </answer>
        </exercise>
        <exercise>
                $\frac{5\pi}{4}$
                <answer>
                    $\frac{\pi}{4}$
                </answer>
        </exercise>
        <exercise>
                $400^{\circ}$
                <answer>
                    $40^{\circ}$
                 </answer>
        </exercise>
        <exercise>
                $\frac{13\pi}{8}$
                <answer>
                    $\frac{3\pi}{8}$
                </answer>
        </exercise>
        <exercise>
                $2$
                <answer>
                    $\pi-2 \approx 1.14$
                </answer>
        </exercise>
        <exercise>
                $\frac{10\pi}{11}$
            <answer>
                $\frac{\pi}{11}$
            </answer>
        </exercise>
        <exercise>
                $2000^{\circ}$
            <answer>
                $20^{\circ}$
            </answer>
        </exercise>
        <exercise>
                $-\frac{9\pi}{5}$
            <answer>
                $\frac{\pi}{5}$
            </answer>
        </exercise>
        <exercise>
                $-100^{\circ}$
            <answer>
                $80^{\circ}$
            </answer>
        </exercise>
    </exercisegroup>
    <exercisegroup cols="2" xml:id="exercisegroup-decimal">
      \subsection*{Convert to Decimal Form}
        
            In <xref ref="exercisegroup-decimal" text="local">Exercises</xref>, convert the angle measure to decimal form (round your answers to the nearest thousandth when necessary).
        
        <exercise>
                $243^{\circ}10'$
            <answer>
               $243^{\circ}10' \approx 243.167^{\circ}$
            </answer>
        </exercise>
        <exercise>
                $3^{\circ}25'$
            <answer>
               $3^{\circ}25' \approx 3.417^{\circ}$
            </answer>
        </exercise>
        <exercise>
                $-23^{\circ}3'$
            <answer>
               $-23^{\circ}3' = -23.05^{\circ}$
            </answer>
        </exercise>
        <exercise>
                $75^{\circ}32'17''$
            <answer>
               $75^{\circ}32'17'' \approx 75.538^{\circ}$
            </answer>
        </exercise>
    </exercisegroup>
    <exercisegroup cols="2" xml:id="exercisegroup-dms">
      \subsection*{Convert to $D^{\circ}M'S''$ Form}
        
            In <xref ref="exercisegroup-dms" text="local">Exercises</xref>, convert the angle measure to $D^{\circ}M'S''$ form.
        
        <exercise>
                $12.4^{\circ}$
            <answer>
               $12.4^{\circ} = 12^{\circ}24'$
            </answer>
        </exercise>
        <exercise>
                $1.53^{\circ}$
            <answer>
               $1.53^{\circ} = 1^{\circ}31'48''$
            </answer>
        </exercise>
        <exercise>
                $-144.9^{\circ}$
            <answer>
               $-144.9^{\circ} = -144^{\circ}54'$
            </answer>
        </exercise>
        <exercise>
                $0.416^{\circ}$
            <answer>
               $0.416^{\circ} = 0^{\circ}24'57.6''$
            </answer>
        </exercise>
    </exercisegroup>
</exercises>
</section>